\bookchapter{Chapitre 67 : Polissage et publication — tests, débogage et mise en ligne sur l'App Store}

\chapterintro{L'App est faite, et ensuite ? Ce chapitre répond aux trois étapes du passage de « ça tourne » à « ça se livre » : comment déboguer avec efficacité pour traquer les bugs, comment écrire des tests pour prévenir les régressions, comment porter l'App sur l'App Store pour que le monde entier la télécharge. La mise en ligne n'est pas qu'une cérémonie — c'est la ligne de partage qui change une réalisation en produit.}

\booksection{Les trois armes du débogage : print, points d'arrêt, aperçu}

\begin{center}
\begin{tabularx}{\linewidth}{|X|X|X|}
\hline
\textbf{Outil} & \textbf{À quoi ça sert} & \textbf{Comment s'en servir} \tabularnewline
\hline
print + console & Lire les journaux d'exécution & Dans le code, print("valeur=\textbackslash{}(x)") ; la sortie se lit dans la console en bas de Xcode \tabularnewline
Débogage par points d'arrêt & Arrêter le programme sur une ligne pour voir les valeurs des variables & Un clic près du numéro de ligne fait naître un marqueur bleu ; l'exécution s'y suspend, et tu avances pas à pas \tabularnewline
Aperçu SwiftUI + débogueur & Voir en direct l'interface et la hiérarchie des vues & Canevas \#Preview + inspecteur View Hierarchy en cours d'exécution \tabularnewline
\hline
\end{tabularx}
\end{center}

\begin{center}\small\color{gray}Les trois armes du débogage\end{center}

\begin{tcbinfo}{Astuce : comment lire un plantage ?}
Quand Swift plante, la console imprime “Fatal error” accompagné de la cause, et Xcode s'arrête à la ligne fautive en la surlignant. L'historique des plantages sur appareil réel se consulte dans Xcode Window Organizer Crashes (compte payant exigé). Cherche d'abord avec les mots-clés « première phrase de la cause de l'erreur + swiftui » — une question précise vaut une réponse précise.

\end{tcbinfo}

\booksection{Les tests : une assurance-vie pour le code}

Le test à la main est traître — tu modifies une fonctionnalité, une autre casse en silence. Les tests valent à la fois « pilule du remords » et « gilet pare-balles ». Sur iOS, le framework XCTest se décline en deux :

\begin{center}
\begin{tabularx}{\linewidth}{|X|X|X|X|}
\hline
\textbf{Niveau} & \textbf{Ce qui est testé} & \textbf{Target} & \textbf{Vitesse} \tabularnewline
\hline
Tests unitaires & Logique pure (calculs, validations) & MyFirstAppTests & En secondes \tabularnewline
Tests UI & Comportement réel de l'interface & MyFirstAppUITests & En minutes \tabularnewline
Test manuel & Ressenti et expérience d'un humain & —— & Toujours nécessaire \tabularnewline
\hline
\end{tabularx}
\end{center}

\begin{center}\small\color{gray}Les niveaux de test iOS (le projet modèle embarque déjà deux Test Target)\end{center}

\begin{lstlisting}[language=Swift]
import XCTest

// extraire en fonction pure la logique de "convertir le texte du montant en nombre"
func parseAmount(_ input: String) -> Double? {
    guard let v = Double(input.trimmingCharacters(in: .whitespaces)),
          v > 0 else { return nil }
    return v
}

final class AmountTests: XCTestCase {
    func testEmptyInputReturnsNil() {
        XCTAssertNil(parseAmount(""))
    }
    func testNegativeReturnsNil() {
        XCTAssertNil(parseAmount("-5"))
    }
    func testNormalValueParses() {
        XCTAssertEqual(parseAmount("12.5"), 12.5)
    }
}
\end{lstlisting}

Lancer les tests : CmdU. Quand tout est au vert, on modifie le code l'esprit tranquille. La validation des montants et les calculs de synthèse du carnet de comptes méritent tous d'être couverts par des tests unitaires.

\booksection{Signature et packaging : apposer le « sceau » sur l'App}

Chaque App iOS doit être signée pour pouvoir s'installer — la preuve que « ce paquet sort bien de toi et n'a pas été altéré ». La hiérarchie :

\begin{itemize}
  \item Phase d'apprentissage : Apple ID gratuit + signature automatique de Xcode (Personal Team) ; ça tourne sur appareil réel, à renouveler tous les 7 jours ;
  \item Phase de distribution : un compte développeur payant (99 dollars/an) délivre certificats et profils, condition de la mise en boutique et de l'installation durable.
\end{itemize}

Packaging : menu Xcode Product Archive ; construction achevée, dans Organizer, Distribute App choisir App Store Connect envoyer. L'archivage doit viser Generic iOS Device comme cible (pas le simulateur) — l'étape que les débutants oublient le plus.

\textbf{Piège à éviter : certificat perdu = mises à jour coupées}\par

Une fois le certificat de signature et la clé privée perdus, il te devient impossible d'envoyer des mises à jour à l'ancienne App — les utilisateurs ne peuvent qu'installer une application « toute neuve ». Sauvegarde soigneusement les certificats (export .p12 du trousseau) et garde le mot de passe en sûreté. Bonne nouvelle : Apple offre la « signature automatique » et la gestion des certificats dans le cloud (case Automatically manage signing dans Xcode), qui épargnent l'essentiel des pièges manuels.

\booksection{Le parcours complet de mise en boutique sur l'App Store}

Créer le compte de développeur

Préparer les éléments de la boutique

Téléverser le build

Remplir les infos de confidentialité et de validation

Validation Apple (1 à 7 jours)

Publier

\begin{center}\small\color{gray}Figure 5 : les six pas de la mise en ligne sur l'App Store\end{center}

\begin{itemize}
  \item Inscription : developer.apple.com, compte individuel 99 dollars/an ;
  \item Éléments : icône (1024×1024, sans coins arrondis — Apple découpe lui-même), captures d'écran (toutes les tailles d'iPhone), description, lien vers la politique de confidentialité ;
  \item Questionnaire de confidentialité : déclare en toute honnêteté quelles données sont recueillies — un carnet de comptes hors ligne peut cocher « ne collecte aucune donnée » ;
  \item Revue : la première soumission attend souvent 1 à 7 jours ; en cas de refus, corrige point par point d'après les « remarques de revue » — les motifs de refus d'Apple sont d'ordinaire très précis, il suffit de les suivre ;
  \item TestFlight : avant la mise en boutique, distribue la version de test à tes amis en bêta interne (jusqu'à 10 000 personnes), récolte un tour de retours puis soumets officiellement — l'allure du vrai professionnel.
\end{itemize}

\booksection{Fabriquer les éléments de boutique : la première impression, c'est le volume de téléchargements}

La fiche boutique est la vitrine que l'utilisateur voit en « flânant » devant le magasin. Trois éléments décident de la première impression :

\begin{itemize}
  \item Icône : 1024×1024, ne te débrouille pas avec une capture d'écran. Sans âme de designer, fais générer un graphique épuré par un outil d'IA ; elle n'est recevable que si, rapetissée sur l'écran d'accueil, on la reconnaît d'un coup d'œil ;
  \item Captures : au moins 2 jeux (aux tailles iPhone dominantes), montrant les scènes fonctionnelles les plus centrales ; chaque capture porte une légende en langue de tous les jours (« une écriture en une seconde » vaut dix fois « une solution de comptabilité performante ») ;
  \item Description : la première phrase établit « ce que c'est + pour qui », sans empiler « le plus pro, le plus fort » — la sincérité rapporte plus que l'enflure.
\end{itemize}

\begin{tcbpractice}{À toi de jouer : « livre » une fois le carnet de comptes}
\booksection{Remontée des plantages et petite liste d'optimisation des performances}

Avant la mise en boutique, ces deux chantiers méritent chacun une demi-heure :

\booksubsection{La surveillance des plantages}

Ce qui n'a pas planté pendant tes tests peut planter sur le téléphone des utilisateurs. Apple livre des rapports de plantage gratuits dans Xcode Organizer (les journaux d'appareils réels supposent un compte payant) ; le choix plus paisible consiste à intégrer un SDK tiers (du genre Firebase Crashlytics). Publier sans surveillance des plantages, c'est rouler les yeux fermés.

\booksubsection{Quatre points d'autocontrôle des performances}

\begin{enumerate}
  \item Les listes avec LazyVStack / List : ne force pas une longue liste dans un simple VStack — le chargement paresseux est la première leçon de performances ;
  \item Les images sans y fourrer les originaux : passe par AsyncImage (image distante) et le cache système ; un original de 5 Mo suffit à étrangler un vieil appareil ;
  \item Pas de gros travaux sur le thread principal : réseau et gros fichiers passent tous par async/await ;
  \item Mesure avec Instruments : l'analyseur de performances embarqué dans Xcode (Time Profiler, détection de fuites mémoire) — mesure d'abord, optimise ensuite, ne te fie pas à ton flair.
\end{enumerate}

(1) Complète les logiques comme parseAmount de 3 tests unitaires (CmdU au vert) ; (2) attrape de tes propres mains, au point d'arrêt, un plantage que tu auras fabriqué, et répare-le ; (3) produis un paquet via Product Archive, puis passe par TestFlight pour inviter des amis en bêta ; (4) si tu as du carburant, inscris-toi au compte développeur et soumets réellement le carnet de comptes en boutique. Les trois premières étapes se bouclent ce soir.

\end{tcbpractice}

\begin{tcbsummary}{Résumé du chapitre}
\begin{itemize}
  \item Trois armes du débogage : journaux print, points d'arrêt, aperçu SwiftUI et inspecteur de vues
  \item XCTest : tests unitaires (CmdU) + tests UI + tests manuels, déjà prévus dans le modèle
  \item Signature : automatique avec un compte gratuit (renouvellement tous les 7 jours) ; seule la distribution exige un compte payant
  \item Packaging : Generic iOS Device + Archive + Distribute ; sauvegarde impérative des certificats
  \item Six pas vers la boutique : inscription éléments envoi questionnaire de confidentialité revue publication ; TestFlight pour la bêta d'abord
\end{itemize}

\end{tcbsummary}
